
        \documentclass[fleqn]{article}      
        \usepackage{amsmath}
        \usepackage{fontspec}
        \usepackage{kotex} % 한국어 지원  

        \begin{document}      
        Sure! Here are the problems divided by difficulty:

 쉬운 문제 (Easy Problems)

\noindent 1. 다음 행렬 \( A = \begin{bmatrix} 1 & 2 \\ 3 & 4 \end{bmatrix} \)의 행과 열의 크기를 구하시오. \\\\

\noindent 2. 두 행렬 \( A = \begin{bmatrix} 2 & 3 \\ 1 & 4 \end{bmatrix} \)와 \( B = \begin{bmatrix} 1 & 0 \\ 2 & 5 \end{bmatrix} \)의 덧셈을 구하시오. \\\\

\noindent 3. 행렬 \( C = \begin{bmatrix} 0 & 5 \\ 2 & 1 \end{bmatrix} \)의 전치 행렬을 구하시오. \\\\

\noindent 4. 행렬 \( A = \begin{bmatrix} 3 & 2 \\ -1 & 4 \end{bmatrix} \)의 행렬식(det)을 구하시오. \\\\

\noindent 5. 행렬 \( A = \begin{bmatrix} 1 & 1 \\ 1 & 1 \end{bmatrix} \)의 스칼라곱을 \( 3 \)로 할 때 결과 행렬을 구하시오. \\\\

 중간 문제 (Medium Problems)

\noindent 6. 행렬 \( A = \begin{bmatrix} 1 & 2 & 3 \\ 4 & 5 & 6 \end{bmatrix} \)와 \( B = \begin{bmatrix} 7 & 8 \\ 9 & 10 \\ 11 & 12 \end{bmatrix} \)의 곱셈 \( AB \)를 구하시오. \\\\

\noindent 7. 행렬 \( A = \begin{bmatrix} 1 & 2 \\ 3 & 4 \end{bmatrix} \)의 역행렬 \( A^{-1} \)을 구하시오. \\\\

\noindent 8. 다음 행렬 \( C = \begin{bmatrix} 1 & 0 & 2 \\ 0 & 1 & 3 \\ 0 & 0 & 1 \end{bmatrix} \)의 고유값을 구하시오. \\\\

\noindent 9. 두 행렬 \( A = \begin{bmatrix} 2 & 3 \\ 5 & 7 \end{bmatrix} \)와 \( B = \begin{bmatrix} 1 & 4 \\ 2 & 6 \end{bmatrix} \)의 행렬식 곱을 구하시오. \\\\

\noindent 10. 다음 행렬 \( A = \begin{bmatrix} 1 & 2 \\ 3 & 5 \end{bmatrix} \)과 \( B = \begin{bmatrix} 4 \\ 5 \end{bmatrix} \)의 곱셈 결과를 구하시오. \\\\

 어려운 문제 (Difficult Problems)

\noindent 11. 다음 행렬 \( A = \begin{bmatrix} 2 & -1 & 0 \\ 0 & 3 & 1 \\ 4 & 0 & 1 \end{bmatrix} \)의 행렬식을 구하시오. \\\\

\noindent 12. 주어진 선형 방정식 \(\begin{bmatrix} 2 & 1 \\ 3 & 4 \\ \end{bmatrix} \begin{bmatrix} x_1 \\ x_2 \end{bmatrix} = \begin{bmatrix} 5 \\ 10 \end{bmatrix}\)의 해를 구하시오. \\\\

\noindent 13. 다음 행렬 \( A = \begin{bmatrix} 1 & 2 \\ -3 & 4 \\ \end{bmatrix} \)와 \( B = \begin{bmatrix} 4 & 0 \\ 2 & 2 \end{bmatrix} \)의 곱의 전치 행렬 \((AB)^T\)을 구하시오. \\\\

\noindent 14. \( 3x + 2y = 6 \)과 \( 4x - y = 5 \)의 계수를 행렬 형태로 표현하고, 이를 행렬을 이용해 \( \begin{bmatrix} x \\ y \end{bmatrix} \)로 나타내시오. \\\\

\noindent 15. 다음 행렬 \( A = \begin{bmatrix} 1 & 2 & 3 \\ 4 & 5 & 6 \\ 7 & 8 & 9 \end{bmatrix} \)의 특성 다항식을 구하시오. \\\\

---

 정답 (Answers)

1. \( 2 \times 2 \) \\\\

2. \( \begin{bmatrix} 3 & 3 \\ 3 & 9 \end{bmatrix} \) \\\\

3. \( \begin{bmatrix} 0 & 2 \\ 5 & 1 \end{bmatrix} \) \\\\

4. \( 10 \) \\\\

5. \( \begin{bmatrix} 3 & 3 \\ 3 & 3 \end{bmatrix} \) \\\\

6. \( \begin{bmatrix} 58 & 64 \\ 79 & 86 \end{bmatrix} \) \\\\

7. \( A^{-1} = \frac{1}{(1)(4) - (2)(3)} \begin{bmatrix} 4 & -2 \\ -3 & 1 \end{bmatrix} = \begin{bmatrix} -4 & 2 \\ 3 & -1 \end{bmatrix} \) \\\\

8. \( 1, 1, 1 \) \\\\

9. \( -1 \) \\\\

10. \( \begin{bmatrix} 30 \\ 35 \end{bmatrix} \) \\\\

11. \( 6 \) \\\\

12. \( x_1 = 2, x_2 = 1 \) \\\\

13. \( \begin{bmatrix} 14 & 6 \\ -6 & -4 \end{bmatrix} \) \\\\

14. \( A \begin{bmatrix} x \\ y \end{bmatrix} = \begin{bmatrix} 6 \\ 5 \end{bmatrix} \) \\\\

15. \( \lambda^3 - 15\lambda^2 + 54\lambda - 0 \) \\\\      
        \end{document} 
    